Glioblastoma (GBM) remains one of the hardest solid tumors to model prognostically because survival is short, recurrence is rapid, and molecular heterogeneity is extreme. Recent glioma studies have shown that outcome is jointly shaped by transcriptomic, methylomic, and genomic structure rather than by any single modality alone \cite{li2025multiomics,piyadasa2026multiomic}. In parallel, single-cell analyses have revealed that tumor ecosystems and cell states contain clinically relevant information that is invisible in bulk measurements \cite{fan2025tme}.

However, three practical gaps remain for GBM-specific survival modeling. First, many published pipelines either focus on mixed glioma cohorts or evaluate pooled cross-validation settings that do not reflect true deployment on a new cohort. Second, the contribution of single-cell priors is often asserted qualitatively, without directly comparing RNA-only, RNA+methylation, and RNA+methylation+scRNA-informed feature sets. Third, deep architectures are frequently introduced without a clear mechanistic link between biological prior knowledge and the actual computation performed by the model.

This study addresses those gaps through a cross-scale benchmark centered on strict CPTAC$\rightarrow$TCGA validation. We curate a multimodal GBM manifest, summarize a 17-patient single-cell atlas into recurrent marker programs, and transfer those programs into both classical and deep survival models. Our benchmark spans LASSO and ridge-style linear baselines, SVM and tree ensembles, a Self-Normalizing Network, a compact transformer, and a prior-guided kernel transformer with biologically masked attention. Rather than claiming that architectural complexity alone solves GBM prognosis, we use this benchmark to identify which modality combinations generalize, which priors remain data-limited, and where the strongest opportunities for methodological improvement still lie.
